\documentclass[12pt]{report}
\usepackage[utf8]{inputenc}
\usepackage[french]{babel}
\usepackage[T1]{fontenc}
\usepackage{amsmath}
\usepackage{amsfonts}
\usepackage{amssymb}
\usepackage{graphicx}
\usepackage{titlesec}
\usepackage{caption}
\usepackage{titling}
\usepackage{booktabs}
\usepackage{enumitem}
\usepackage{eurosym}
\usepackage{epigraph}
\usepackage{hyperref}
\usepackage{fontspec}
\usepackage{ragged2e}
\usepackage{parskip}
\usepackage{wrapfig}
\usepackage{calc}
\usepackage{float}

\graphicspath{ {../../static/} {img/} }
\setlength{\droptitle}{-10em}
\titleformat{\chapter}[hang]{\normalfont\huge\bfseries}{\thechapter. }{0em}{}

\begin{document}

\title{
	{\vspace{3em}\protect\centering\protect\includegraphics[width=0.9\textwidth]{pacification_vector.pdf}}\\
	{\vspace{4em}\Huge Rapport de soutenance}\\
	{\large Brainless Devs}
}
\author{
	Thibault Allançon\\
	Valérian Fayt
	\and
	Antoine Gonzalez\\
	Cédric Parpet}
\date{
	{\vfill\protect\centering\protect\includegraphics{brainless_devs.pdf}}\\
	Dossier Projet Informatique\\
	Info-Sup EPITA\\
	Mai 2018
}

\maketitle
\tableofcontents

\chapter{Introduction}

La deuxième phase de développement de \textit{Pacification} touche à sa fin. Durant ces quelques semaines, nous avons pu beaucoup progresser sur l’évolution du jeu, qui commence à réellement prendre forme. 

Nous somme partis d’un éditeur de map avec des fonctionnalités basiques pour arriver à un solo et un multijoueur jouables, avec des unités capables d'interagir avec leur environnement. Le réseau transmet correctement les informations, la map dispose d’une génération procédurale et d’un brouillard de guerre, et tous nos assets sont prêts.

Malgré l’apparition de nouvelles difficultés, tous les membres du groupe ont pu travailler sur leurs parties et le résultat de cette période est conforme à nos attentes.

Dans ce rapport, nous expliquerons avec plus de précision tout ce qui a été réalisé durant cette phase, ainsi que ce qu’il reste à faire avant la soutenance finale.

\begin{figure}
    \centering
    \includegraphics[width=0.6\textwidth]{project_mood_S2}
    \caption*{\textit{Calvin and Hobbes}, Bill Watterson}
\end{figure}

\chapter{Cahier des charges}

Cette fois-ci, la répartition des tâches nous convient parfaitement. 

Les assets étant finalement terminés, l’interface - actuellement en retard, plutôt 70\% que 80\% - pourra avancer rapidement et être terminée sous peu.

Le site web n’a pas réellement eu besoin d’être amélioré depuis la première soutenance, mis à part une mise à jour des informations.

La map est presque terminée, la majorité de ses fonctionnalités sont implémentées, il faut maintenant attendre une progression des autres parties afin de la compléter. On peut donc estimer l’avancement de la map à 95\% au lieu de 100 (qui ne pourra réellement être atteint qu’à la fin).

L’IA a progressé comme prévue, son fonctionnement sera détaillé plus bas.

Le gameplay a également atteint ses objectifs, les unités sont implémentées et fonctionnelles, les villes sont utilisables, le Player fait partie intégrante du jeu, il ne reste maintenant qu’à implémenter une économie au jeu grâce aux ressources et aux bâtiments des villes ; et à prendre les biomes en compte dans les calculs des unités.

\chapter{Avancement}

\section{Map (Thibault)}

La première chose qu'il a fallu intégrer à la map était les unités en elles mêmes. Le système de pathfinding était déjà en place, mais le fait de positionner les unités, de pouvoir les sélectionner, les déplacer, représentait une tâche importante et majeure du jeu. 

En utilisant des courbes de Bézier nous avons désormais des déplacements fluides et assez naturels sur la carte, d'autant plus que ces courbes nous permettent de gérer facilement l'orientation des unités qui se tournent avant et pendant leurs déplacements.

\begin{figure}[H]
    \centering
    \includegraphics[width=0.8\textwidth]{404}
    \caption{Unités et courbes de bézier}
\end{figure}

Lors de la soutenance précédente nous avions mis en place des menus pour éditer la carte afin de montrer les différentes fonctionnalités. Nous avons décidé de réutiliser cette partie pour mettre en place un éditeur de map afin de créer ses propres cartes sur lesquelles on pourra ensuite jouer.

\begin{figure}[H]
    \centering
    \includegraphics[width=0.8\textwidth]{404}
    \caption{Editeur de map}
\end{figure}

Cependant, il était prévu depuis le début de générer des cartes lors d'un début de partie pour avoir à chaque nouveau lancement de jeu une situation différente.

Une génération de carte procédurale est désormais intégrée au jeu, et le joueur peut maintenant sélectionner s'il souhaite utiliser une carte qu'il a lui même édité ou bien en générer une nouvelle.

\begin{figure}[H]
    \centering
    \includegraphics[width=0.8\textwidth]{404}
    \caption{Génération de carte}
\end{figure}

La dernière fonctionnalité que nous voulions intégrer à la carte de Pacification était le brouillard de guerre, élément indispensable pour ce type de jeu de stratégie. Il a fallu adapter les shaders ainsi que les textures pour cacher ou montrer des parties de la carte en fonction de plusieurs critères comme : la hauteur de la case sur laquelle se trouve l'unité, son nombre de point de vision, si la case a déjà été visité ou non. On différencie donc deux types de brouillard : d'une part les terrains inexplorés qui sont totalement cachés du joueur et de l'autre les terrain précédemments explorés mais hors de vision et portée de l'unité.

\begin{figure}[H]
    \centering
    \includegraphics[width=0.8\textwidth]{404}
    \caption{Brouillard de guerre}
\end{figure}

\section{Réseau et multijoueur (Valérian)}

\begin{figure}[H]
    \centering
    \includegraphics[width=0.8\textwidth]{404}
    \caption{TODO}
\end{figure}

\section{Gameplay (Antoine)}

Le gameplay a bien avancé. Nous disposons à présent d’unités correctement implémentées.

Au début de la partie, le joueur commence avec un colon et un fantassin. Le colon permet de construire la première ville du joueur, tandis que l’unité offensive peut permettre d’explorer les alentours et de découvrir la carte, alors recouverte par le brouillard de guerre. 

\begin{figure}[H]
    \centering
    \includegraphics[width=0.8\textwidth]{404}
    \caption{Situation initiale}
\end{figure}

Les unités peuvent donc se déplacer (pour le moment sans animation développée, elles glissent sur la map et pivotent) et activer leur action propre (construire une ville, exploiter une ressource, attaquer une unité ennemie).

\begin{figure}[H]
    \centering
    \includegraphics[width=0.8\textwidth]{404}
    \caption{Pathfinding}
\end{figure}

La ville permet pour le moment de produire d’autres unités grâce au menu “barrack”. Il est ainsi possible de créer un autre colon, un travailleur, ou une unité offensive, qui apparaîtront sur la ville.

Chaque unité possède des stats qui lui sont propres, ainsi qu’un modèle 3D assigné à sa création selon son type (les unités d’attaque possèdent un second modèle amélioré qu’elles obtiennent à partir du niveau 11). Ces statistiques, présentées dans les tableaux de la première soutenance, influeront notamment sur le nombre de tour nécessaire pour tuer une cible et  la vitesse de déplacement.

La ville disposera également d’améliorations, cependant elles dépendent de la population, et l’économie sera implémentée lors de la dernière phase de développement.

\section{Assets (Cédric)}

Les Assets ont été finis à 100\%, elles sont basées, pour la plupart, sur mes premières créations pour faire les suivantes. Nous sommes donc restés sur le style graphique sur lequel nous avons commencé. Trois nouvelles textures ont été créées  spécialement pour les villes.

\begin{figure}[H]
    \centering
    \includegraphics[width=0.8\textwidth]{404}
    \caption{BlueRoof, RedRoof, DarkStoneWall}
\end{figure}

Les modèles 3D des bâtiments sont tous finis et et à l’exception des villes aucune n’a demandé de nouveau modèle à faire. Les derniers bâtiments pour les ressources ont été modélisés et comme pour les premières ressources déjà incluses dans le jeu, il existe deux modèles différents : la version naturelle et la version exploitée 

\begin{figure}[H]
    \centering
    \includegraphics[width=0.8\textwidth]{404}
    \caption{foret naturelle et exploitée}
\end{figure}

Concernant les villes la première version a été légèrement modifiée, et les nouvelles textures créées permettent une meilleure différenciation des différents modèles correspondant aux différentes ères.

\begin{figure}[H]
    \centering
    \includegraphics[width=0.8\textwidth]{404}
    \caption{villes}
\end{figure}

Concernant les unités ajoutées il y a les deux versions d’unité à portée, les deux versions des unités de siège (catapulte et canon), et le cavalier qui correspond à l’amélioration de l’infanterie. Comme pour les modèles précédent, la création des unités (notamment le canon et la catapulte) se sont fait à l’aide d’image de référence. Les modèles des unités humaines ont été construits à partir du modèle humain de base créer durant la première période du projet. Toutes les unités ont été agrandies pour faciliter leur sélection. Certains bug graphiques ont été corrigés.

Les animations des personnages ont été considéré comme du bonus. Cependant des recherches ont quand même été faites sur cet aspect là : La catapulte a eu une première animation d’attaque, de déplacement et de destruction qui pourront pour certaines être réutilisé pour le canon. Des essais ont été faits sur l’animation de l’armature des unités humaines cependant celle-ci s’est trouvé très compliqué à gérer. Cependant le site internet Mixamo.com permet de donner une armature à un personnage humanoïde que l’on a créé et propose un panel d’animations très fluides et gratuites.

\section{Interface (Cédric, Valérian, Antoine)}

Pour le menu nous avons récupérer un menu créé par MICHSKY qui l’a rendu accessible à tous gratuitement sur internet. Nous l’avons aménagé de manière à ce qu’il s’adapte parfaitement à notre jeu.

\begin{figure}[H]
    \centering
    \includegraphics[width=0.8\textwidth]{404}
    \caption{MainMenu}
\end{figure}

 Le menu principal comporte un accès au “menu Play” qui regroupe les fonctionnalités du jeu en solo et du jeu en multijoueur.

\begin{figure}[H]
    \centering
    \includegraphics[width=0.8\textwidth]{404}
    \caption{MenuPlay}
\end{figure}

 Les images apparaissant sur les boutons “Solo” “Host” et “Join” ont été faites grâce à blender. Les positions du soldat ont été possible grâce à l’armature et aux animations importés sur le site Mixamo.com cité prècedemment. Un menu “Paramètres” a été créé où l’on pourra par exemple changer la langue du jeu ou activer/désactiver du jeu. Un bouton “information”s est disponible et mène aux crédits du jeu.  On navigue de menu en menu grâce à des déplacements de caméras qui nous amènent sur des nouveaux canvas et des boutons back ont été placé pour simplifier cette navigation.

%Le logo pacification a été créé sur le site FontStruct.com en s’inspirant d’un modèle trouvé sur internet. Pour insérer le titre correctement importer la police d’écriture ne suffisait pas car le texte apparaissait flou. Nous avons donc mis une image du texte. Pour cela nous avons utilisé le logiciel open source Krita qui a permis d’écrire et de rendre le fond transparent afin que le mot “Pacification” apparaisse sans fond sur le menu.

De la musique a été ajoutée dans le jeu.

L’interface de la caserne (barrack) est reliée aux villes. Sa position et son échelle actuelle seront peut-être amenée à changer, l’interface actuelle n’étant pas encore parfaitement organisée. 

\section{IA (Thibault)}

\begin{figure}[H]
    \centering
    \includegraphics[width=0.8\textwidth]{404}
    \caption{TODO}
\end{figure}

\section{Site Web (Valérian)}

\begin{figure}[H]
    \centering
    \includegraphics[width=0.8\textwidth]{404}
    \caption{TODO}
\end{figure}

\chapter{Avances, retards, difficultés rencontrées}

\section*{Map}

\begin{figure}[H]
    \centering
    \includegraphics[width=0.8\textwidth]{404}
    \caption{TODO}
\end{figure}

\section*{Réseau et Multijoueur}

\begin{figure}[H]
    \centering
    \includegraphics[width=0.8\textwidth]{404}
    \caption{TODO}
\end{figure}

\section*{Gameplay}

Si on reprend la hiérarchie des classes présentée lors de la première soutenance, on peut voir que les unités sont complètement implémentées, et les villes le sont en partie (l’économie reste à implémenter). Le “Player” est également géré par le réseau et permet de coordonner les villes et les unités. Les biomes sont présents dans le code, mais pour le moment ils ne sont pas pris en compte par le pathfinding et la génération de la map: il reste à organiser une probabilité d’apparition des différentes ressources.

\begin{figure}[H]
    \centering
    \includegraphics[width=1\textwidth]{class_hierarchy}
    \caption{hiérarchie (avancement)}
\end{figure}

Pour la soutenance finale, il reste donc à implémenter l’économie du jeu (donc les bâtiments des villes et les ressources) et adapter les paramètres des biomes/troupes afin d’obtenir un jeu le plus équilibré possible.

La difficulté principale durant cette période a été d'apprendre à "lier" les classes selon Unity (qui doivent être instanciées pour être affichées sur la map), et celles de C\# qui appellent un constructeur (dont j'avais besoin pour gérer les unités et les villes), sans pour autant créer de conflits avec les autres systèmes eux aussi nécessaires au fonctionnement du jeu. Cela m'a pris plus de temps que prévu et a été une source de stress, mais tout a fini par fonctionner.

\section*{Assets}

\begin{figure}[H]
    \centering
    \includegraphics[width=0.8\textwidth]{404}
    \caption{TODO}
\end{figure}

\section*{Interface}

\begin{figure}[H]
    \centering
    \includegraphics[width=0.8\textwidth]{404}
    \caption{TODO}
\end{figure}

\section*{IA}

\begin{figure}[H]
    \centering
    \includegraphics[width=0.8\textwidth]{404}
    \caption{TODO}
\end{figure}

\section*{Site web}

\begin{figure}[H]
    \centering
    \includegraphics[width=0.8\textwidth]{404}
    \caption{TODO}
\end{figure}

\chapter{À venir}

Pour la troisième et dernière soutenance, nous allons devoir terminer \textit{Pacification}. Ce qui signifie mettre en place l'économie et équilibrer le gameplay, finir l'interface pour la rendre la plus confortable et intuitive possible, ajouter les ressources à la map, terminer l'IA, et bien sûr prendre le tout en compte dans le réseau qu'il faudra stabiliser.

En somme, beaucoup de travail, mais rien que nous ne puissions faire à temps. Le jeu ayant déjà bien avancé, et les systèmes principaux étant déjà mis en place, nous saurons ajouter ces dernières fonctionnalités sans difficulté majeure.

\chapter{Expériences personnelles}

\begin{itemize}
	\item \textbf{Thibault :} TODO
        \item \textbf{Valérian :} TODO
	\item \textbf{Cédric :} TODO
	\item \textbf{Antoine :} Pour ma part, je suis satisfait de mes avancées, bien que, comme expliqué plus haut, je n'avais pas anticipé de telles difficultés pour lier les systèmes entre eux et les faire fonctionner correctement. Cependant, j'ai pu apprendre de ces problèmes et je sais maintenant les résoudre, ce qui me facilitera la tâche lors de la dernière phase de développement du projet.
\end{itemize}

\chapter{Conclusion}

\begin{figure}[H]
    \centering
    \includegraphics[width=0.8\textwidth]{404}
    \caption{TODO}
\end{figure}

\end{document}
